\section{Il modello di Machine Learning}
\label{sec:modello-ml}

\subsection{Origine, dataset e scelta architetturale}

Il modello di riconoscimento utilizzato in PomodorIA non è stato addestrato
da zero nell'ambito di questo progetto, ma riutilizza il lavoro svolto in un
esame precedente~[\ref{bib:cnntomato}], in cui è stata confrontata
sperimentalmente una rete completamente connessa (FCNN) con una rete
convoluzionale (CNN) sul dataset \textbf{PlantVillage -- Tomato}~[\ref{bib:dataset}]
(14.529 immagini, 10 classi: 9 patologie fogliari più la classe "sana").

La Tabella~\ref{tab:fcnn-vs-cnn} riassume l'esito di quel confronto, che ha
motivato la scelta della CNN come modello da portare in produzione nel
sistema Edge: una rete completamente connessa tratta l'immagine come un
vettore piatto di pixel, perdendo l'informazione sulla struttura spaziale, e
per questo confonde più facilmente patologie visivamente simili; la CNN,
grazie ai filtri convoluzionali, cattura pattern spaziali locali (bordi,
texture, forma delle macchie fogliari) risultando nettamente più accurata.

\begin{table}[h]
\centering
\small
\begin{tabular}{@{}p{4.6cm}p{4.2cm}p{4.2cm}@{}}
\toprule
& \textbf{FCNN [2L-512-ReLU]} & \textbf{CNN [64f-k3-3blk]} \\
\midrule
Accuracy (holdout 80/20) & $\sim$84\% & $\sim$96\% \\
Recall minimo per classe & 0,69 (Target Spot) & 0,90 \\
Punto debole & Confonde patologie fungine simili & Classi rare più critiche \\
\bottomrule
\end{tabular}
\caption{Confronto tra rete completamente connessa (FCNN) e rete convoluzionale (CNN).}
\label{tab:fcnn-vs-cnn}
\end{table}

\subsection{Architettura della CNN: \texttt{TomatoCNN}}

Il modello scelto per il deployment Edge è la configurazione
\texttt{CNN~[64f-k3-3blk]}: 64 filtri iniziali, kernel $3\times3$, 3 blocchi
convoluzionali, per un totale di 8.765.066 parametri (33,44 MB su disco in
formato float32). L'architettura è composta da:

\begin{itemize}
    \item una sequenza di \textbf{3 blocchi convoluzionali} (\texttt{Conv2d}
    $\to$ \texttt{ReLU} $\to$ \texttt{MaxPool2d}), che raddoppiano il numero
    di filtri (64 $\to$ 128 $\to$ 256) e dimezzano la dimensione spaziale
    dell'immagine ad ogni blocco ($64\times64 \to 8\times8$);
    \item un \textbf{classificatore} completamente connesso
    (\texttt{Linear} $\to$ \texttt{ReLU} $\to$ \texttt{Dropout} $\to$
    \texttt{Linear}) che produce 10 valori grezzi (\emph{logits}), uno per
    classe -- la loro trasformazione in probabilità (softmax) è applicata
    dal motore di inferenza, descritto nella Sezione~\ref{sec:edge-layer}.
\end{itemize}

\subsection{Preprocessing dell'immagine}

Ogni immagine, prima di essere passata alla rete, viene ridimensionata a
$64\times64$ pixel RGB e normalizzata con media e deviazione standard
tipiche del dataset ImageNet ($\text{mean}=[0.485, 0.456, 0.406]$,
$\text{std}=[0.229, 0.224, 0.225]$) -- la stessa pipeline usata in fase di
addestramento, condizione necessaria perché la rete funzioni correttamente
in inferenza.
